\documentclass{article}

% Language and encoding
\usepackage[portuguese]{babel}
\usepackage[utf8]{inputenc}
\usepackage[T1]{fontenc}

% Page layout
\usepackage[a4paper,left=25mm,right=25mm,top=25mm,bottom=25mm,headheight=6mm,footskip=12mm]{geometry}
\setlength{\parindent}{1em}
\usepackage{booktabs}
\usepackage{indentfirst}
\usepackage{ragged2e}
\usepackage{setspace}
\usepackage{caption}
\usepackage{float}
\usepackage{longtable}
\usepackage{enumitem}
\usepackage{graphicx}
\usepackage{adjustbox}
\usepackage{lmodern}
\usepackage{xcolor}
\usepackage{listings}
\usepackage{lastpage}
\usepackage{fancyhdr}
\usepackage{hyperref}

%% headers & footers
\fancyhf{}
\rfoot{\small{\thepage\ / \pageref{LastPage}}}
\pagestyle{fancy}
\renewcommand{\headrulewidth}{0.4pt}
\renewcommand{\footrulewidth}{0.4pt}

%% colors
\definecolor{engineering}{rgb}{0.549,0.176,0.098}
\definecolor{cloudwhite}{cmyk}{0,0,0,0.025}

%% source-code listings
\lstset{
 language=bash,
 basicstyle=\footnotesize\ttfamily,
 keywordstyle=\bfseries,
 numbers=left,
 numberstyle=\scriptsize\texttt,
 stepnumber=1,
 numbersep=8pt,
 frame=tb,
 float=htb,
 aboveskip=8mm,
 belowskip=4mm,
 backgroundcolor=\color{cloudwhite},
 showstringspaces=false,
 showtabs=false,
 tabsize=2,
 captionpos=b,
 belowcaptionskip=12pt,
 breaklines=true,
 breakatwhitespace=false,
 commentstyle=\color{gray},
 lineskip=3pt,
 extendedchars=true
}

%% hyperreferences
\hypersetup{
 plainpages=false,
 pdfpagelayout=SinglePage,
 bookmarksopen=false,
 bookmarksnumbered=true,
 breaklinks=true,
 linktocpage,
 colorlinks=true,
 linkcolor=engineering,
 urlcolor=engineering,
 filecolor=engineering,
 citecolor=engineering,
 allcolors=engineering
}

%% path to figures
\graphicspath{{figures/}{figuras/}}

%% macros
\newcommand{\school}{Universidade do Minho}
\newcommand{\degree}{Licenciatura em Engenharia Informática}
\newcommand{\projtitle}{Projeto de Processamento de Linguagens}
\newcommand{\subtitle}{2025/2026}
\newcommand{\projauthor}{Trabalho Prático}
\newcommand{\projname}{Relatório Técnico - Compilador Fortran 77 para EWVM}

\newcommand{\atsfigure}[3]{
\begin{figure}[H]
\centering
\IfFileExists{#1}{
    \includegraphics[width=#2\textwidth]{#1}
}{
    \fbox{\parbox[c][50mm][c]{0.9\textwidth}{\centering Placeholder para print: #1}}
}
\caption{#3}
\end{figure}
}

\renewcommand{\thesection}{\arabic{section}}
\setcounter{secnumdepth}{3}

\begin{document}
\justifying

% Capa
\begin{titlepage}
\centering

\vspace{-15mm}
{\large \textbf{\textsc{\school}}}\\

\vspace{0.5cm}
{\large \degree}\\[8mm]

\vspace{2cm}
\IfFileExists{UM.png}{\includegraphics[width=52mm]{UM.png}}{\vspace{20mm}}

\vspace{1cm}
{\Large \projtitle}\\[2mm]
{\large \subtitle}\\[20mm]

{\Large \textbf{\projauthor}}\\
\vspace{0.5cm}
{\Large \projname}

\vspace{4cm}

{\large Francisco Barros A106894}\\[2mm]
{\large Tiago Martins A106927}\\[2mm]
{\large Beatriz Miranda A107365}\\[2mm]

\vspace{2cm}
{\Large \textbf{Maio 2026}}

\end{titlepage}

\renewcommand{\contentsname}{Índice}
\tableofcontents
\newpage
\listoffigures
\newpage
\listoftables
\newpage

% ─────────────────────────────────────────────────────────────────────────────
\section{Introdução}
% ─────────────────────────────────────────────────────────────────────────────

Este relatório tem como objetivo apresentar o trabalho prático desenvolvido na unidade
curricular de Processamento de Linguagens. O trabalho consiste no desenvolvimento de um
compilador para um \textit{subset} de Fortran 77, tendo como destino final a máquina virtual
disponibilizada (EWVM). Deste modo, apresentamos a arquitetura do sistema, as decisões de
implementação tomadas pelo grupo, as dificuldades encontradas e os resultados obtidos.

Para além da implementação base, o projeto inclui funcionalidades de valorização, como suporte
a subprogramas externos (\texttt{FUNCTION}, \texttt{SUBROUTINE}, \texttt{CALL} e
\texttt{RETURN}), tratados por \textit{inlining} na geração de IR, e tratamento robusto de
arrays multidimensionais com indexação dinâmica. Na iteração final foram ainda reforçados o
pré-processamento de entrada, a compatibilidade com programas em formato livre simples, o
\textit{backend} de I/O, a validação semântica de labels e ciclos, a linearização correta de
arrays segundo a ordem de Fortran e a cobertura dos testes \textit{end-to-end}. O objetivo
deste documento é apresentar, de forma estruturada, todas as fases da \textit{pipeline} de
compilação, bem como as opções arquiteturais que sustentam a entrega.

% ─────────────────────────────────────────────────────────────────────────────
\section{Arquitetura}
% ─────────────────────────────────────────────────────────────────────────────

A implementação foi organizada como uma \textit{pipeline} de compilação com fronteiras claras entre fases, de forma a separar responsabilidades e reduzir acoplamento entre componentes. Em vez de concentrar toda a lógica num único módulo, cada etapa transforma uma representação de entrada numa representação de saída mais estruturada, permitindo testar e validar o compilador de forma incremental.

Em termos práticos, o fluxo segue a sequência:
\begin{enumerate}
\item normalização do código Fortran 77 em \textit{fixed-form} ou formato livre simples;
\item tokenização do código normalizado;
\item construção da AST;
\item validação semântica;
\item geração de representação intermédia;
\item otimização;
\item emissão de código EWVM.
\end{enumerate}

A principal vantagem desta arquitetura é que cada módulo pode evoluir de forma relativamente independente. Por exemplo, melhorias no parser não obrigam a alterar o pré-processamento, desde que o contrato entre as duas fases seja preservado. Do mesmo modo, a fase de otimização opera sobre IR e não depende diretamente de detalhes da sintaxe concreta da linguagem fonte.

Outra decisão importante foi preservar informação contextual útil (como número de linha e labels) ao longo das primeiras fases. Esta informação é essencial para a semântica (por exemplo, validação de \texttt{GOTO}) e para mensagens de erro mais informativas.

\atsfigure{pipeline.png}{0.40}{Pipeline do compilador (visão geral).}

% ─────────────────────────────────────────────────────────────────────────────
\subsection{Pré-processador}
% ─────────────────────────────────────────────────────────────────────────────

O primeiro componente executado é o \textbf{preprocessor.py}, responsável por normalizar o
texto fonte antes da análise léxica. O compilador aceita o formato clássico
\textit{fixed-form} de Fortran 77 e, adicionalmente, um formato livre simples, compatível
com programas copiados diretamente do enunciado. Esta opção reduz o risco de rejeitar
exemplos corretos apenas por diferenças de alinhamento, sem abandonar a semântica de labels
característica de Fortran 77.

A normalização segue as regras tradicionais:
\begin{itemize}
\item colunas 1--5: campo de label;
\item coluna 6: marca de continuação;
\item colunas 7--72: código efetivo da instrução.
\end{itemize}

Com base nestas regras, e nas extensões aceites para formato livre, o módulo executa quatro
tarefas essenciais:
\begin{enumerate}
\item remove linhas de comentário e linhas vazias;
\item identifica labels e separa-as do conteúdo da instrução;
\item junta corretamente linhas de continuação quando a coluna 6 indica continuação;
\item remove comentários de fim de linha, preservando corretamente texto entre aspas.
\end{enumerate}

O resultado não é apenas texto limpo: cada linha lógica passa a ser representada com metadados
(número de linha original, label e conteúdo normalizado). Esta escolha facilita as fases
seguintes, que podem operar sobre uma estrutura uniforme e já livre de ruído lexical.

A manutenção da interpretação \textit{fixed-form} é importante porque preserva a convenção
histórica de colunas e a deteção posicional de labels de \texttt{DO}/\texttt{CONTINUE}.
O suporte a formato livre simples é uma camada de conveniência controlada: permite linhas
como \texttt{PROGRAM T}, \texttt{20 IF (...) THEN} ou \texttt{PRINT *, X}, mas continua a
produzir a mesma representação lógica usada pelo resto da \textit{pipeline}.

Em resumo, o pré-processador funciona como uma fase de estabilização da entrada: reduz ambiguidades precoces e garante que o lexer recebe uma sequência previsível de linhas lógicas.

% ─────────────────────────────────────────────────────────────────────────────
\subsection{Lexer}
% ─────────────────────────────────────────────────────────────────────────────

A segunda fase é implementada em \textbf{lexer.py}, usando \texttt{ply.lex}. O objetivo é converter cada linha lógica normalizada numa sequência de tokens tipados, já preparados para análise sintática.

A implementação define explicitamente:
\begin{itemize}
\item palavras reservadas da linguagem (\texttt{PROGRAM}, \texttt{IF}, \texttt{DO}, \texttt{READ}, \texttt{CALL}, etc.);
\item literais (\texttt{INT\_LIT}, \texttt{REAL\_LIT}, \texttt{BOOL\_LIT}, \texttt{STRING\_LIT});
\item operadores aritméticos, relacionais e lógicos;
\item separadores e símbolos estruturais.
\end{itemize}

Do ponto de vista de robustez, duas decisões foram relevantes:
\begin{enumerate}
\item \textbf{case-insensitive}: o lexer é construído para aceitar diferentes combinações de maiúsculas/minúsculas, refletindo o uso típico de Fortran;
\item \textbf{normalização de identificadores}: os identificadores são normalizados internamente para maiúsculas, garantindo consistência nas fases seguintes.
\end{enumerate}

A tokenização é feita por linha lógica e preserva o contexto de origem (linha e label). Assim, o parser não recebe apenas tokens isolados, mas sim uma sequência contextualizada de \texttt{(lineno, label, token\_stream)}.

Em termos de cobertura da linguagem, esta fase já resolve vários pontos críticos:
\begin{itemize}
\item operadores pontuados de Fortran (\texttt{.EQ.}, \texttt{.NE.}, \texttt{.AND.}, \texttt{.OR.}, \texttt{.NOT.});
\item literais reais com notação científica, incluindo expoente \texttt{E} e \texttt{D};
\item strings delimitadas por aspas simples ou duplas, incluindo apóstrofos duplicados no estilo Fortran;
\item formas alternativas de palavras-chave relevantes, como \texttt{GO TO};
\item distinção entre palavras reservadas e identificadores válidos.
\end{itemize}

Quando surge um caractere inválido, o lexer sinaliza erro imediatamente, evitando propagar estados inconsistentes para o parser. Esta estratégia reduz o custo de depuração e melhora a previsibilidade da pipeline.

% ─────────────────────────────────────────────────────────────────────────────
\subsection{Parser}
% ─────────────────────────────────────────────────────────────────────────────

O \textbf{parser.py} segue uma abordagem híbrida, combinando \texttt{ply.yacc} com análise estrutural por linhas. Esta decisão foi tomada para conciliar dois objetivos: cumprir o requisito de utilização de PLY e, ao mesmo tempo, lidar bem com construções clássicas de Fortran 77 orientadas a label.

A parte suportada diretamente por \texttt{ply.yacc} cobre:
\begin{itemize}
\item expressões aritméticas, relacionais e lógicas com precedência/associatividade;
\item instruções de linha única (declarações, atribuições, \texttt{PRINT}, \texttt{READ},
\texttt{WRITE}, \texttt{GOTO}, \texttt{GO TO}, \texttt{CALL}, \texttt{RETURN}, \texttt{STOP});
\item declarações \texttt{CHARACTER} simples e com tamanho explícito, como
\texttt{CHARACTER*12}.
\end{itemize}

Já as estruturas de bloco multi-linha são tratadas por controlo estrutural explícito:
\begin{itemize}
\item \texttt{IF (...) THEN ... ELSE ... ENDIF}, incluindo a grafia \texttt{END IF};
\item \texttt{DO label var = start, end [, step] ... label CONTINUE}.
\end{itemize}

Esta separação traz vantagens práticas:
\begin{enumerate}
\item mantém a gramática LALR mais simples e focada nas construções locais;
\item evita complexidade artificial para padrões de bloco dependentes de labels;
\item melhora a qualidade das mensagens de erro em estruturas incompletas (por exemplo, \texttt{IF} sem \texttt{ENDIF} ou \texttt{DO} sem fecho).
\end{enumerate}

Outro aspeto importante é a desambiguação entre acesso a array e chamada de função em formas
do tipo \texttt{ID(...)}. O parser faz esta decisão com base no contexto conhecido de
subprogramas e funções intrínsecas, produzindo nós AST distintos (\texttt{ArrayRefNode} ou
\texttt{FuncCallNode}) conforme o caso. Nesta fase é também suportada a forma \texttt{F()},
permitindo funções sem argumentos.

Para evitar ruído numa utilização normal do compilador, a construção dos parsers PLY é feita
com \texttt{NullLogger}. Assim, avisos internos sobre tokens que pertencem a outra gramática
parcial não aparecem no terminal durante a compilação de exemplos, mantendo o CLI limpo para
demonstração.

No final da fase sintática, o resultado é uma AST completa, com corpo principal e subprogramas externos, pronta para validação semântica. Assim, o parser funciona como ponte entre uma entrada ainda textual (tokens por linha) e uma representação estrutural rica, adequada às fases de análise e geração de código.

% ─────────────────────────────────────────────────────────────────────────────
\subsection{Análise Semântica}
% ─────────────────────────────────────────────────────────────────────────────

O módulo \textbf{semantic.py} implementa as validações necessárias para garantir a
corretude do programa antes da geração de código. Esta fase atua sobre a AST e usa uma
tabela de símbolos por escopos, separando o programa principal dos subprogramas
(\texttt{FUNCTION} e \texttt{SUBROUTINE}).

A análise é feita de forma estruturada: primeiro são registadas assinaturas globais dos
subprogramas (nome, tipo de retorno, parâmetros), e depois cada corpo é validado no seu
escopo próprio. Esta estratégia evita ambiguidades em chamadas e permite validar referências
cruzadas com antecedência.

Durante o percurso da AST, são verificadas as seguintes propriedades:

\begin{itemize}
\item \textbf{declaração e uso de identificadores}: cada variável, array ou subprograma deve
estar declarado antes de ser usado no respetivo escopo;
\item \textbf{compatibilidade de tipos}: atribuições e operações são validadas para evitar
combinações inválidas; quando aplicável, são aceites promoções seguras (por exemplo, de
inteiro para real);
\item \textbf{validação de operadores}: operadores aritméticos exigem operandos numéricos;
\texttt{.AND.}, \texttt{.OR.} e \texttt{.NOT.} exigem operandos lógicos; comparações
ordenadas exigem operandos numéricos; \texttt{**} é aceite apenas com expoente inteiro; esta
validação impede combinações inválidas como \texttt{LOGICAL + INTEGER} ou
\texttt{CHARACTER * INTEGER};
\item \textbf{condições lógicas}: expressões de controlo em \texttt{IF} têm de produzir
valor lógico;
\item \textbf{labels e desvio}: cada \texttt{GOTO} é validado contra labels realmente
existentes no bloco analisado e labels duplicadas são rejeitadas;
\item \textbf{regras de ciclos DO}: variável de controlo, limites e passo são verificados com
as restrições semânticas da linguagem, incluindo rejeição de passo literal zero;
\item \textbf{arrays}: são validadas dimensões, rank, tipo dos índices e acessos fora da forma
esperada;
\item \textbf{subprogramas}: número, tipo e categoria dos argumentos (escalar/array) devem
coincidir com a assinatura de \texttt{FUNCTION}/\texttt{SUBROUTINE}.
\end{itemize}

Em caso de erro, o compilador interrompe a pipeline nesta fase, evitando que um programa
semanticamente inválido avance para IR ou \textit{backend}. Esta opção melhora a depuração,
uma vez que concentra a deteção de incoerências na fase em que a informação semântica ainda
é mais rica.

% ─────────────────────────────────────────────────────────────────────────────
\subsection{Geração de IR}
% ─────────────────────────────────────────────────────────────────────────────

\paragraph{Representação intermédia usada}
A linguagem intermédia usada nesta fase é \textbf{TAC} (\textit{Three-Address Code}).
Só depois da otimização em TAC é gerado o código final para a \textbf{EWVM}.
\\

O módulo \textbf{ir\_gen.py} é responsável pelo \textit{lowering} da AST para uma
representação intermédia em TAC (\textit{Three-Address Code}). Esta representação desacopla
o \textit{front-end} do compilador do \textit{backend}, facilitando a otimização e a
posterior geração de código VM.

Cada instrução TAC é modelada com operador e operandos (\texttt{op}, \texttt{result},
\texttt{arg1}, \texttt{arg2}), o que torna explícita a dependência entre cálculos e
controlo de fluxo. Ao longo desta fase são gerados temporários e labels para linearizar
estruturas de alto nível em passos elementares.

A tradução cobre, entre outros, os seguintes casos:

\begin{itemize}
\item \textbf{expressões}: operadores aritméticos, relacionais e lógicos são convertidos em
sequências TAC com temporários;
\item \textbf{atribuições e I/O}: leitura, escrita e cópia de valores são normalizadas em
operações simples;
\item \textbf{controlo de fluxo}: \texttt{IF/ELSE} e \texttt{DO} são convertidos para labels e
saltos condicionais/incondicionais;
\item \textbf{arrays}: acessos e atualizações tornam-se operações explícitas de carga e
armazenamento indexado;
\item \textbf{subprogramas}: chamadas são baixadas para IR por inlining, com preservação da
semântica através de mapeamento de símbolos, etiquetas internas, renomeação sistemática de
variáveis locais e saltos de retorno.
\end{itemize}

No caso de \texttt{PRINT}/\texttt{WRITE} list-directed, a IR inclui separadores explícitos
entre valores, permitindo que uma instrução como \texttt{PRINT *, A, B} produza uma saída
legível equivalente a \texttt{A B}, em vez de concatenar diretamente os valores. Nos ciclos
\texttt{DO}, a geração de IR emite também a label real associada ao \texttt{CONTINUE}
terminal, garantindo que um \texttt{GOTO} para essa label tem destino válido no código VM.

Ao final desta etapa, o compilador deixa de depender da forma sintática original do programa:
toda a informação relevante passa a estar codificada numa sequência TAC uniforme, pronta para
otimização e geração de código final.

\atsfigure{ir_exemplo.png}{0.40}{Excerto de IR (TAC) gerado.}

% ─────────────────────────────────────────────────────────────────────────────
\subsection{Otimizador}
% ─────────────────────────────────────────────────────────────────────────────

O módulo \textbf{optimizer.py} aplica um conjunto de \textit{passes} de otimização
sobre o IR antes da geração de código final para a VM. Em vez de uma única
transformação, o sistema executa uma sequência de passes até atingir \textbf{ponto
fixo}, isto é, até duas iterações consecutivas produzirem exatamente o mesmo IR.

\paragraph{Passes implementados}
Atualmente, o otimizador inclui os seguintes passes principais:

\begin{enumerate}
\item \textbf{Constant folding}: avalia operações binárias com operandos numéricos
constantes (por exemplo, \texttt{ADD 2 3}) e substitui-as por uma cópia do valor
resultante.
\item \textbf{Copy propagation}: propaga aliases introduzidos por atribuições
simples (\texttt{COPY}), reescrevendo usos posteriores para reduzir temporários
intermédios.
\item \textbf{Dead temporary elimination}: remove definições de temporários
(\texttt{\_t*}) cujo valor nunca é utilizado, desde que a instrução seja segura
de eliminar.
\item \textbf{Eliminação de código inalcançável}: remove instruções após saltos
incondicionais não seguidos de label, que nunca seriam executadas.
\item \textbf{Eliminação de escritas mortas}: remove atribuições a variáveis cujo
valor é sobrescrito antes de ser lido.
\item \textbf{Peephole}: aplica simplificações locais, nomeadamente remoção de
cópias redundantes (\texttt{x = x}) e saltos incondicionais para a label
imediatamente seguinte.
\end{enumerate}

\paragraph{Observação sobre alcance}
No estado atual, estas são as otimizações do módulo de otimização propriamente dito.
Existem ainda simplificações noutras fases do compilador (por exemplo, \textit{inlining}
de subprogramas no lowering para IR e resolução direta de alguns acessos a arrays
com índices constantes no backend), mas não fazem parte dos passes de
\texttt{optimizer.py}.


% ─────────────────────────────────────────────────────────────────────────────
\subsection{Geração de Código VM}
% ─────────────────────────────────────────────────────────────────────────────

O módulo \textbf{codegen.py} traduz o IR (TAC) para as instruções da EWVM. O
\textit{backend} implementa as seguintes funcionalidades:

\begin{itemize}
  \item mapeamento de variáveis para memória global;
  \item linearização de arrays multidimensionais em ordem \textit{column-major}, coerente
  com Fortran;
  \item suporte a índices constantes e dinâmicos;
  \item I/O tipado no código gerado, com separadores em escrita list-directed;
  \item emissão segura de literais string no código VM;
  \item tradução do operador de exponenciação \texttt{**} para um ciclo VM de
  multiplicações sucessivas (instrução intermédia \texttt{POW}).
\end{itemize}

O operador \texttt{**} merece nota especial. Em vez de depender de uma instrução nativa
de potência na EWVM, o backend emite um ciclo que inicializa o resultado a 1 e multiplica
pela base enquanto o expoente inteiro for positivo. Esta solução mantém a tradução portável
dentro do conjunto de instruções já usado pelo projeto.

Para arrays, o backend distingue dois casos. Quando todos os índices são literais inteiros,
o deslocamento é resolvido estaticamente e o código usa diretamente \texttt{PUSHG}/
\texttt{STOREG}. Quando algum índice é dinâmico, o deslocamento linear é calculado em tempo
de execução e são emitidas instruções de acesso indireto (\texttt{LOADN}/\texttt{STOREN}).
Esta combinação melhora a eficiência nos casos constantes e preserva a generalidade nos casos
com variáveis de índice.

\atsfigure{vm_exemplo.png}{0.40}{Excerto de código VM gerado.}

% ─────────────────────────────────────────────────────────────────────────────
\section{Implementação}
% ─────────────────────────────────────────────────────────────────────────────

Concluída a visão geral da arquitetura, este capítulo aprofunda as decisões de implementação
mais relevantes, analisando de que modo as diferentes funcionalidades do \textit{subset} de
Fortran 77 foram tratadas ao longo da pipeline.

% ─────────────────────────────────────────────────────────────────────────────
\subsection{Controlo de fluxo: IF e DO}
% ─────────────────────────────────────────────────────────────────────────────

Uma das principais dificuldades do Fortran 77 em \textit{fixed-form} é a gestão de blocos
com labels. As construções IF/ENDIF e DO/.../CONTINUE não seguem uma estrutura de blocos
delimitados por chavetas ou palavras-chave simples, mas sim por labels numéricas associadas
a instruções CONTINUE.

A solução adotada foi o \textit{parser} estrutural por linhas tokenizadas, que percorre a
sequência de linhas e identifica os delimitadores de bloco com base nas labels. A validação
semântica complementa esta análise, garantindo a coerência das labels (e.g., ausência de
GOTO para labels inexistentes, estruturas DO sem CONTINUE correspondente).

Para a geração de IR, cada bloco de controlo é emitido com labels únicas, evitando colisões
em programas com múltiplos ciclos ou condicionais. A versão final emite ainda o rótulo VM
correspondente ao \texttt{CONTINUE} terminal do ciclo \texttt{DO}, uma vez que Fortran permite
saltos explícitos para esse label. Sem esta emissão, programas sintaticamente válidos com
\texttt{GOTO} para o fim do ciclo poderiam falhar apenas na execução da VM.

% ─────────────────────────────────────────────────────────────────────────────
\subsection{Declarações e tipos}
% ─────────────────────────────────────────────────────────────────────────────

As declarações de variáveis são processadas durante a análise semântica, que constrói uma
tabela de símbolos por âmbito (programa principal e cada subprograma). A compatibilidade de
tipos é verificada em atribuições e expressões aritméticas, com emissão de erros descritivos
para facilitar a depuração por parte do utilizador.

A versão atual reforça ainda a validação ao nível dos operadores: cada operador exige
operandos do tipo correto, o que impede combinações semanticamente inválidas mesmo que a
sintaxe esteja bem formada. Foram também acrescentadas validações para labels duplicadas e
para \texttt{DO} com passo literal zero. Esta validação aproxima o compilador do contrato
semântico esperado de uma linguagem tipada.

% ─────────────────────────────────────────────────────────────────────────────
\subsection{Arrays multidimensionais}
% ─────────────────────────────────────────────────────────────────────────────

O suporte a arrays multidimensionais implicou duas decisões de implementação relevantes.
Em termos de semântica, as dimensões declaradas são validadas contra os acessos no código.
Em termos de \textit{backend}, os arrays são linearizados por \textit{stride}: o endereço
de cada elemento é calculado em tempo de compilação para índices constantes, e emitido como
sequência de instruções aritméticas para índices dinâmicos.

A linearização segue a ordem \textit{column-major} de Fortran, isto é, o primeiro índice é o
que varia mais rapidamente. Por exemplo, em \texttt{A(2,3)}, os elementos \texttt{A(1,2)} e
\texttt{A(2,1)} têm deslocamentos diferentes dos que seriam obtidos numa ordem típica de C.
Esta correção é essencial para compatibilidade semântica com Fortran 77.

% ─────────────────────────────────────────────────────────────────────────────
\subsection{I/O: READ, PRINT e WRITE}
% ─────────────────────────────────────────────────────────────────────────────

As instruções \texttt{READ}, \texttt{PRINT} e \texttt{WRITE} são traduzidas para as primitivas
de I/O da EWVM. O \textit{backend} seleciona a instrução adequada consoante o tipo da expressão
(inteiro, real, lógico ou string), garantindo que o código gerado realiza a leitura e escrita
com o formato correto. A escrita list-directed introduz separadores entre valores, evitando
saídas ambíguas como \texttt{12} quando o programa pretende escrever \texttt{1} e \texttt{2}
como campos distintos.

% ─────────────────────────────────────────────────────────────────────────────
\subsection{Subprogramas: FUNCTION e SUBROUTINE}
% ─────────────────────────────────────────────────────────────────────────────

O suporte a subprogramas externos constitui a principal funcionalidade de valorização do
projeto. A implementação introduz desafios em três fases distintas.

Na análise semântica, são validadas as assinaturas e os argumentos de cada chamada, assim
como o tipo de retorno das FUNCTIONs. No \textit{lowering} para IR, cada subprograma é
inlinado com um mapeamento interno de símbolos e labels, evitando colisões de nomes em
múltiplas chamadas.

Na versão final, o inlining foi corrigido para isolar corretamente variáveis locais.
Anteriormente, apenas parâmetros e valores de retorno recebiam nomes internos únicos; uma
variável local com o mesmo nome de uma variável do programa principal podia causar colisão
silenciosa. Agora, as declarações locais do subprograma são renomeadas com prefixos únicos
por chamada e emitidas no IR através de instruções declarativas. Assim, preserva-se a noção
de âmbito mesmo sem implementar uma \textit{stack} de chamadas nativa na EWVM. O
comando \texttt{RETURN} é materializado como salto para um rótulo de retorno, ficando o
backend encarregado apenas de emitir a sequência TAC já linearizada.

Esta decisão é semanticamente importante: a transformação por inlining só é correta se o
código resultante for equivalente ao programa original. A renomeação sistemática de símbolos
locais garante precisamente essa equivalência.

% ─────────────────────────────────────────────────────────────────────────────
\subsection{Exemplos representativos}
% ─────────────────────────────────────────────────────────────────────────────

Tal como no relatório de referência, a profundidade técnica beneficia de exemplos concretos que
mostram o percurso entre o código fonte, a AST, o IR e o código final para a EWVM. Os casos que
se seguem são representativos das principais decisões de implementação do compilador.

\paragraph{Programa elemental: \texttt{HELLO}}
O exemplo mais simples valida o funcionamento de ponta a ponta da pipeline sem interferência de
controlo de fluxo ou estruturas compostas.

\begin{verbatim}
      PROGRAM HELLO
      PRINT *, 'Ola, Mundo!'
      END
\end{verbatim}

O código gerado é compacto e directo, confirmando a tradução base da instrução \texttt{PRINT}:

\begin{verbatim}
PUSHN 0
START
PUSHS "Ola, Mundo!"
WRITES
WRITELN
STOP
\end{verbatim}

\paragraph{Controlo de fluxo: \texttt{FATORIAL}}
O programa \texttt{FATORIAL} demonstra a tradução de ciclos \texttt{DO} com label, bem como a
linearização do corpo do ciclo em TAC e na EWVM.

\begin{verbatim}
      PROGRAM FATORIAL
      INTEGER N, I, FAT
      READ *, N
      FAT = 1
      DO 10 I = 1, N
      FAT = FAT * I
 10   CONTINUE
      PRINT *, FAT
      END
\end{verbatim}

Neste caso, o backend materializa o ciclo com labels dedicadas, teste de continuação e salto de
retorno ao início do ciclo, garantindo um comportamento equivalente ao esperado em Fortran 77.

\paragraph{Arrays multidimensionais: \texttt{SOMAARR}}
O programa \texttt{SOMAARR} é útil para demonstrar indexação dinâmica em arrays e o cálculo de
deslocamentos efectivos em tempo de execução.

\begin{verbatim}
      PROGRAM SOMAARR
      INTEGER NUMS(5)
      INTEGER I, SOMA
      SOMA = 0
      DO 30 I = 1, 5
      READ *, NUMS(I)
      SOMA = SOMA + NUMS(I)
 30   CONTINUE
      PRINT *, SOMA
      END
\end{verbatim}

Quando o índice é uma variável, o backend já não pode resolver o endereço de forma estática,
recorrendo por isso a uma sequência explícita de instruções para calcular o deslocamento.

\paragraph{Subprogramas: \texttt{CONVERSOR}}
O exemplo \texttt{CONVERSOR} ilustra a integração entre o programa principal e uma função
externa, bem como o uso de \texttt{MOD}, \texttt{IF} e \texttt{GOTO} no interior do subprograma.

\begin{verbatim}
      PROGRAM CONVERSOR
      INTEGER NUM, BASE, RESULT, CONVRT
      READ *, NUM
      DO 10 BASE = 2, 9
      RESULT = CONVRT(NUM, BASE)
      PRINT *, RESULT
 10   CONTINUE
      END

      INTEGER FUNCTION CONVRT(N, B)
      INTEGER N, B, QUOT, REM, POT, VAL
      VAL = 0
      POT = 1
      QUOT = N
 20   IF (QUOT .GT. 0) THEN
      REM = MOD(QUOT, B)
      VAL = VAL + (REM * POT)
      QUOT = QUOT / B
      POT = POT * 10
      GOTO 20
      ENDIF
      CONVRT = VAL
      RETURN
      END
\end{verbatim}

Este caso é particularmente relevante porque confirma a estratégia de \textit{inlining} adotada para
subprogramas. A função é validada semanticamente, convertida para IR com etiquetas de retorno
e símbolos locais renomeados e, no final, integrada no fluxo principal sem necessidade de uma
convenção de chamada nativa na EWVM.

% ─────────────────────────────────────────────────────────────────────────────
\section{Testes e Validação}
% ─────────────────────────────────────────────────────────────────────────────

A validação do compilador inclui testes unitários e de pipeline que cobrem todas as fases da
compilação. Os testes estão organizados pelas seguintes categorias:

\begin{itemize}
  \item preprocessamento;
  \item análise léxica;
  \item análise sintática;
  \item semântica (casos positivos e negativos);
  \item geração de IR;
  \item otimizador;
  \item \textit{backend} VM;
  \item arrays multidimensionais;
  \item subprogramas;
  \item exemplos oficiais do enunciado.
\end{itemize}

Estado de execução no ambiente de desenvolvimento após a iteração final:
\textbf{92 passed}.

Os testes de integração com a EWVM, que anteriormente poderiam depender de ligação à máquina
de referência, correm localmente graças ao módulo \textbf{vm.py}, um interpretador Python do
subconjunto de instruções EWVM gerado pelo compilador. Esta solução elimina a dependência de
rede nos testes automáticos, mantendo a cobertura \textit{end-to-end} completa.

Para além disso, o módulo \textbf{optimize.py} disponibiliza uma interface de linha de comandos
para otimizar representações IR textuais independentemente do compilador, útil para depuração
e validação isolada do otimizador.

Foram acrescentados testes específicos para os pontos de maior risco de desvalorização:
programas em formato livre simples, linhas com indentação leve, \texttt{END IF},
\texttt{READ(*,*)}, \texttt{WRITE(*,*)}, separação de valores em \texttt{PRINT},
labels duplicadas, passo zero em \texttt{DO} e ordem \textit{column-major} em arrays
multidimensionais.


\atsfigure{testes_pytest.png}{0.40}{Resultado da suite de testes no ambiente local.}

% ─────────────────────────────────────────────────────────────────────────────
\section{Compatibilidade com o enunciado}
% ─────────────────────────────────────────────────────────────────────────────

Os exemplos oficiais do enunciado foram usados como referência funcional e como base de
validação da pipeline completa. A implementação cobre os programas de teste fornecidos e os
casos de maior exigência prática, incluindo ciclos \texttt{DO}, condicionais aninhadas,
acessos a arrays multidimensionais e chamadas a subprogramas com passagem de parâmetros.

Do ponto de vista da conformidade técnica, o compilador respeita os requisitos centrais do
trabalho: utiliza \texttt{ply.lex} no lexer, \texttt{ply.yacc} no parser, uma fase semântica
separada e uma representação intermédia em TAC antes da geração de código para a EWVM. A
estratégia adotada para \texttt{IF} e \texttt{DO} privilegia a clareza estrutural, ao passo que
\texttt{FUNCTION} e \texttt{SUBROUTINE} são tratados por \textit{inlining} com isolamento de
variáveis locais, o que simplifica o \textit{backend} sem comprometer a semântica observável.

A opção por combinar \texttt{ply.yacc} com análise estrutural por linhas não é a solução mais
convencional, mas revelou-se a mais adequada ao carácter de Fortran 77, sobretudo em programas
com labels e blocos dependentes de contexto. Uma gramática puramente LALR para todas estas
formas tornaria o código mais frágil e as mensagens de erro menos claras.

Em síntese, a implementação foi desenhada para cumprir o enunciado com rigor, mas também para
produzir um compilador tecnicamente legível, modular e defensável em prova oral.

% ─────────────────────────────────────────────────────────────────────────────
\section{Avaliação pelos critérios do enunciado}
% ─────────────────────────────────────────────────────────────────────────────

\subsection{Correção}

O compilador processa corretamente os cinco exemplos padrão do enunciado:
\texttt{HELLO}, \texttt{FATORIAL}, \texttt{PRIMO}, \texttt{SOMAARR} e
\texttt{CONVERSOR}. Estes exemplos exercitam a pipeline completa, incluindo I/O, ciclos
\texttt{DO}, \texttt{IF/ELSE}, \texttt{GOTO}, arrays e \texttt{FUNCTION}. Os ficheiros VM
correspondentes foram regenerados a partir do compilador atual e os testes
\textit{end-to-end} confirmam os resultados esperados no interpretador local da EWVM.

Além dos exemplos oficiais, existem testes negativos para garantir que programas incorretos
são rejeitados na fase adequada: identificadores não declarados, tipos incompatíveis, labels
inexistentes ou duplicadas, dimensões inválidas de arrays, índices fora de limites literais e
ciclos \texttt{DO} com passo inválido. Assim, a correção não se limita a aceitar programas
válidos; inclui também recusar programas semanticamente incoerentes.

\subsection{Estrutura}

O código está organizado por fases clássicas de compilação: \textbf{preprocessor.py},
\textbf{lexer.py}, \textbf{parser.py}, \textbf{semantic.py}, \textbf{ir\_gen.py},
\textbf{optimizer.py} e \textbf{codegen.py}. Esta separação reduz acoplamento e torna cada
responsabilidade verificável de forma isolada. A AST é definida em \textbf{ast\_nodes.py}, a
IR em \textbf{ir.py} e a execução local de validação em \textbf{vm.py}. A estrutura é, por
isso, adequada tanto para manutenção como para defesa, pois permite explicar claramente onde
cada decisão técnica é aplicada.

\subsection{Funcionalidade}

A implementação cobre as construções essenciais solicitadas: declarações de tipos e
variáveis, expressões aritméticas, relacionais e lógicas, \texttt{IF-THEN-ELSE}, ciclos
\texttt{DO} com labels, \texttt{GOTO}/\texttt{GO TO}, \texttt{READ}, \texttt{PRINT},
\texttt{WRITE}, arrays, \texttt{STOP} e subprogramas externos. Para valorização, suporta
\texttt{FUNCTION}, \texttt{SUBROUTINE}, \texttt{CALL}, \texttt{RETURN}, funções sem
argumentos, arrays multidimensionais com índices dinâmicos e exponenciação \texttt{**}.
Também aceita variantes frequentes em Fortran, como \texttt{END IF}, \texttt{READ(*,*)},
\texttt{WRITE(*,*)}, \texttt{CHARACTER*n} e literais reais com expoente \texttt{D}.

\subsection{Eficiência}

A eficiência é tratada em dois níveis. Primeiro, há uma etapa explícita de otimização sobre
TAC, com \textit{constant folding}, \textit{copy propagation}, eliminação de temporários
mortos, eliminação de código inalcançável, eliminação de escritas mortas e simplificações
\textit{peephole}. Segundo, o \textit{backend} evita trabalho desnecessário quando possível:
índices constantes de arrays são resolvidos estaticamente, enquanto índices dinâmicos usam
endereçamento indireto apenas quando necessário. O código VM gerado é, portanto, mais simples
do que uma tradução totalmente direta e satisfaz o requisito de estar minimamente otimizado.

\subsection{Defesa}

Para a defesa, o projeto é demonstrável por comandos simples:
\texttt{python3 src/main.py examples/hello.f -o out.vm}, execução dos testes com
\texttt{python3 -m pytest -q} e inspeção opcional de \texttt{--dump-ast} e
\texttt{--dump-ir}. A equipa deve estar preparada para explicar a razão da arquitetura em
fases, a utilização de PLY, a escolha de TAC, a estratégia de \textit{inlining} para
subprogramas, a linearização de arrays em ordem Fortran e os passes de otimização. Estes são
os pontos que demonstram conhecimento técnico individual para além da simples execução do
programa.

% ─────────────────────────────────────────────────────────────────────────────
\section{Trabalho Futuro}
% ─────────────────────────────────────────────────────────────────────────────

Após a conclusão do projeto, identificaram-se algumas áreas onde o compilador poderia ser
melhorado em trabalho futuro:

\begin{itemize}
  \item maior cobertura de integração com a EWVM em ambiente de integração contínua (CI);
  \item expansão da bateria de programas longos de referência para validação de robustez;
  \item implementação de uma convenção de chamadas real na VM, substituindo a abordagem por
        inlining e permitindo suporte a recursão;
  \item reforço de validações semânticas em casos fronteira (e.g., recursão em
        subprogramas, \textit{aliasing} de arrays);
  \item extensão da semântica para mais intrínsecas de Fortran 77;
  \item suporte adicional a mais variantes sintáticas de Fortran 77 não cobertas pelo
        subconjunto atual;
  \item exploração de otimizações adicionais no IR, como eliminação de subexpressões
        comuns e \textit{loop-invariant code motion}.
\end{itemize}

% ─────────────────────────────────────────────────────────────────────────────
\section{Conclusão}
% ─────────────────────────────────────────────────────────────────────────────

Concluindo este projeto com sucesso, o grupo considera que o desenvolvimento do mesmo
constituiu uma excelente oportunidade para aplicar na prática os conceitos de processamento
de linguagens, desde a análise léxica até à geração de código para uma máquina virtual real.

O projeto apresenta uma implementação modular e consistente de um compilador Fortran 77
para EWVM, com cobertura dos requisitos base e dos principais pontos de valorização. A
pipeline completa, o desenho por fases e a validação automatizada sustentam uma entrega
tecnicamente sólida e uma defesa bem fundamentada. Em particular, na versão final foram
reforçados o pré-processamento, o suporte a variantes sintáticas usuais, o backend de I/O,
a validação semântica de labels e ciclos, a linearização de arrays multidimensionais segundo
a ordem de Fortran e a cobertura dos testes \textit{end-to-end}.

A estrutura adotada --- separação clara entre \textit{front-end} (preprocessamento, lexer,
parser e semântica), camada intermédia (IR/TAC e otimização) e \textit{backend} (geração
VM) --- melhora a manutenibilidade, facilita a depuração por fases e permite introduzir
melhorias de forma incremental sem comprometer a estabilidade global do compilador.

Os resultados obtidos na suite de testes (92 passed) demonstram robustez nos casos
essenciais do enunciado e nos cenários de valorização, incluindo subprogramas e arrays com
indexação dinâmica. Em contexto de defesa, isto permite sustentar não apenas que o compilador
``funciona'', mas também que as decisões de implementação foram tecnicamente justificadas,
testáveis e alinhadas com critérios de correção, estrutura e eficiência.

% ─────────────────────────────────────────────────────────────────────────────
\newpage
\appendix
\renewcommand{\thesection}{\Alph{section}}
% ─────────────────────────────────────────────────────────────────────────────

\section{Gramática}
\label{anexo:gramatica}

\begin{verbatim}
program        := "PROGRAM" ID stmt* "END" subprogram*
subprogram     := function_def | subroutine_def
function_def   := type "FUNCTION" ID "(" [id_list] ")" stmt* "END"
subroutine_def := "SUBROUTINE" ID "(" [id_list] ")" stmt* "END"

stmt           := decl
               | assign | print | read
               | if_stmt | do_stmt | goto
               | continue | call | return | stop

decl           := type var_spec ("," var_spec)*
type           := "INTEGER" | "REAL" | "LOGICAL" | "CHARACTER"
               | "CHARACTER" "*" INT_LIT
var_spec       := ID | ID "(" expr_list ")"
assign         := lvalue "=" expr
lvalue         := ID | ID "(" expr_list ")"
print          := "PRINT" "*" "," expr_list
               | "WRITE" "*" "," expr_list
               | "WRITE" "(" "*" "," "*" ")" expr_list
read           := "READ" "*" "," lvalue_list
               | "READ" "(" "*" "," "*" ")" lvalue_list
if_stmt        := "IF" "(" expr ")" "THEN" stmt* ["ELSE" stmt*] ("ENDIF" | "END" "IF")
do_stmt        := "DO" label ID "=" expr "," expr ["," expr]
                  stmt* label "CONTINUE"
goto           := "GOTO" (INT_LIT | ID)
               | "GO" "TO" (INT_LIT | ID)
call           := "CALL" ID "(" [expr_list] ")"

expr           := literal | ID | ID "(" [expr_list] ")"
               | "(" expr ")" | "-" expr | ".NOT." expr
               | expr ("+"|"-"|"*"|"/"|"**") expr
               | expr (".EQ."|".NE."|".LT."|".LE."|".GT."|".GE.") expr
               | expr (".AND."|".OR.") expr
\end{verbatim}

% ─────────────────────────────────────────────────────────────────────────────
\section{Matriz de conformidade com o enunciado}
\label{anexo:conformidade}

\begin{table}[H]
\centering
\begin{tabular}{p{0.46\textwidth} p{0.16\textwidth} p{0.28\textwidth}}
\toprule
\textbf{Requisito} & \textbf{Estado} & \textbf{Evidência} \\
\midrule
Lexer com PLY                   & Cumprido & lexer.py \\
Parser com PLY Yacc             & Cumprido & parser.py \\
Análise semântica               & Cumprido & semantic.py \\
Validação de tipos e operadores & Cumprido & semantic.py \\
IF/DO/GOTO                      & Cumprido & parser + IR + testes \\
READ/PRINT/WRITE                & Cumprido & parser/codegen + testes \\
Exponenciação (**)              & Cumprido & codegen (ciclo POW) \\
Subprogramas (valor.)           & Cumprido & parser/semantic/IR \\
Isolamento de variáveis locais  & Cumprido & ir\_gen.py \\
Funções sem argumentos          & Cumprido & parser.py \\
Formato livre simples           & Cumprido & preprocessor.py \\
END IF / GO TO                  & Cumprido & parser.py \\
Arrays em ordem Fortran         & Cumprido & codegen.py \\
Otimização IR                   & Cumprido & optimizer.py \\
Testes exemplos oficiais EWVM   & Cumprido & test\_ewvm\_outputs.py \\
\bottomrule
\end{tabular}
\caption{Resumo de conformidade técnica.}
\label{tab:compliance}
\end{table}

% ─────────────────────────────────────────────────────────────────────────────
\section{Instruções de execução}
\label{anexo:execucao}

\subsection{Preparação}
\begin{verbatim}
python3 -m venv .venv
source .venv/bin/activate
pip install -r requirements.txt
\end{verbatim}

\subsection{Compilação}
\begin{verbatim}
python3 src/main.py examples/hello.f -o out.vm
python3 src/main.py examples/fatorial.f --dump-ast --dump-ir -o out.vm
\end{verbatim}

\subsection{Testes}
\begin{verbatim}
python -m pytest -q
RUN_EWVM_TESTS=1 python -m pytest -q tests/test_ewvm_outputs.py
\end{verbatim}

\end{document}
